\subsection{Dimensión y profundidad de un Módulo normalmente plano a lo largo de un
subpreesquema cerrado.}
\label{subsection:IV.6.10}

\begin{env}[6.10.1]
\EGAReaderTarget{ega4.c04.statement.6-10-1}{IV.2.s6.10.1-fr}
Sean $X$ un preesquema localmente noetheriano, $\mathcal J$ un Ideal
cuasicoherente de $\mathcal O_X$, $Y$ el subpreesquema cerrado de $X$ definido
por $\mathcal J$, $j:Y\to X$ la inyección canónica. Para todo entero $k\geq0$,
el $\mathcal O_X$-Módulo
$\mathcal J^k\mathcal F/\mathcal J^{k+1}\mathcal F$ es anulado por
$\mathcal J$ y, por tanto, tiene la forma $j_*(\mathcal G_k)$, donde
$\mathcal G_k=j^*(\mathcal J^k\mathcal F/\mathcal J^{k+1}\mathcal F)=j^*(\mathcal J^k\mathcal F)$
es un $\mathcal O_Y$-Módulo coherente. Por abuso de lenguaje, denotaremos por
$\operatorname{gr}_{\mathcal J}^{\bullet}(\mathcal F)$ el
$\mathcal O_Y$-\emph{Módulo graduado} igual a la suma directa
\[
\bigoplus_{k=0}^{\infty}\mathcal G_k
=j^*\left(\bigoplus_{k=0}^{\infty}
\mathcal J^k\mathcal F/\mathcal J^{k+1}\mathcal F\right);
\]
en particular, se tiene
$\operatorname{gr}_{\mathcal J}^{0}(\mathcal F)=\mathcal G_0
=\mathcal F\otimes_{\mathcal O_Y}\mathcal O_Y=j^*(\mathcal F)$. Diremos
(siguiendo a Hironaka) que $\mathcal F$ es \emph{normalmente plano a lo largo de} $Y$ si
$\operatorname{gr}_{\mathcal J}^{\bullet}(\mathcal F)$ es un
$\mathcal O_Y$-\emph{Módulo plano}. Es equivalente ((2.1.12) y
$(\mathbf{0}_{\mathrm I},6.1.2)$) a decir que cada uno de los $\mathcal O_Y$-Módulos
$\mathcal G_k=j^*(\mathcal J^k\mathcal F)$ es \emph{localmente libre}.
\end{env}

\begin{proposition}[6.10.2]
\label{IV.2.s6.10.2-fr}
Sean $X$ un preesquema localmente noetheriano, $Y$ un subpreesquema cerrado
\textup{íntegro} de $X$. Para todo $\mathcal O_X$-Módulo coherente $\mathcal F$,
existe un abierto $U$ de $X$ tal que $Y\cap U\ne\varnothing$ y tal que
$\mathcal F|U$ es normalmente plano a lo largo de $Y\cap U$.
\end{proposition}

\begin{proof}
Sea, en efecto, $\mathcal J$ el Ideal coherente de $\mathcal O_X$ que define
$Y$; la $\mathcal O_Y$-Álgebra
$\mathcal B=\operatorname{gr}_{\mathcal J}^{\bullet}(\mathcal O_X)$ es
cuasicoherente y \emph{de tipo finito}, pues está generada por
$\operatorname{gr}_{\mathcal J}^{1}(\mathcal O_X)$, imagen inversa de
$\mathcal J/\mathcal J^2$; como
$\operatorname{gr}_{\mathcal J}^{\bullet}(\mathcal F)$ es un
$\mathcal B$-módulo cuasicoherente generado por
$\operatorname{gr}_{\mathcal J}^{1}(\mathcal F)$, es un
$\mathcal B$-Módulo \emph{de tipo finito}. Si se pone
$Z=\operatorname{Spec}(\mathcal B)$, el morfismo estructural $u:Z\to Y$ es
entonces de tipo finito, y si $\mathcal H$ es el $\mathcal O_Z$-Módulo coherente
tal que $u_*(\mathcal H)=\operatorname{gr}_{\mathcal J}^{\bullet}(\mathcal F)$,
basta aplicar a $u$ y a $\mathcal H$ el teorema de planitud
genérica (6.9.1) para demostrar la proposición.
\end{proof}
